\section{The Split}

The geometry-to-gauge conversion creates a new danger. Because the gauge is
generated by boundary geometry, the reader may try to carry the geometric
sector and the gauge sector as two independent summands. The grammar forbids
that finite direct sum. A finite certificate may read the geometric face or
the gauge face, but it may not form an interior mixed scalar that pretends
both faces are simultaneously available as independent components.

The split is therefore a carrier discipline. It says how a reading may pass
through the Hilbert door without losing the distinction between the
geometric source and the gauge sign. The geometric carrier reads boundary
separation, curvature-like residue, and the mesh on which the ledger is
refined. The gauge carrier reads transported orientation, charge, and the
discriminating action. The same invariant underlies both, but the finite
reader is not allowed to add them as if it possessed a larger interior space
where their mixed product were already meaningful.

The reason is the mixed term. If geometry and gauge were freely summed, an
interior scalar cross-reading would be available. That cross-reading would
say that the geometric residue and the gauge residue can be multiplied
inside the same finite certificate without remainder. But the chapter has
not earned such a certificate. The boundary generated the gauge precisely
because the interior could not settle the charged residue on its own. To
reinsert the gauge as an interior summand would erase the boundary
obligation that generated it.

Orthogonality is the disciplined alternative. It does not mean the two
sectors are unrelated. It means the finite grammar reads them in separated
channels whose interior scalar mixing is forbidden. The geometric channel
can carry the boundary-to-interior refinement. The gauge channel can carry
the sign transport. What cannot be carried is a hidden interior scalar that
lets one channel cancel debt in the other. The split protects the charged
reading from being dissolved into geometry and protects geometry from being
reduced to a gauge label.

This explains why boundary radiation is waiting in the next chapter. If the
mixed term cannot live in the interior, then any reconciliation between the
geometric and gauge faces must appear at the boundary. Radiation will be the
grammar's name for that boundary phenomenon. The present chapter does not
need to derive it yet. It only has to force the condition that makes it
unavoidable: no interior mixed certificate.

The split also clarifies matter and antimatter readings. Matter and
antimatter are not two substances thrown into a common box. They are signs
read through the gauge under specified baselines. The native flat baseline
reads the electron side. A tilted baseline reads the positron side. If the
geometric and gauge sectors were collapsed into a direct sum, the baseline
dependence would blur. The grammar would no longer know whether a positive
sign came from a tilt, a boundary, or an illicit interior mixture. The split
keeps the sign honest.

The Yang--Mills illustration helps at the level of labels. Different
refinement classes may require local repair rules that close under their own
composition and commute only where their labels are independent. The smooth
shadow may look like a product of symmetries, but the ledger discipline is
more severe: independent label classes must not corrupt one another's
bookkeeping. The present split is even sharper. Geometry and gauge are not
two free label classes to be added in the interior. They are two faces of one
boundary-forced reading, and their mixed scalar is unavailable.

The split is not a failure of unification. It is the form unification takes
under finite measurement. A loose unification would announce one large space
and place everything inside it. The grammar refuses. It identifies the
single invariant underneath both readings, but it also preserves the carrier
discipline that makes each reading meaningful. The unity is in the invariant
and the boundary process. The split is in the finite certificates by which
the readings are carried.

Thus the geometry/gauge relation is neither separation without contact nor
fusion without discipline. Geometry generates the gauge at the boundary.
Gauge reads the charged residue. The finite Hilbert pairing door then forces
a split carrier: geometric or gauge, not an illicit sum. This is why the
chapter can say that the sectors do not combine as a direct sum while still
saying that one grammar has forced both faces.

